% Part: first-order-logic
% Chapter: syntax-and-semantics
% Section: free-vars-sentences

\documentclass[../../../include/open-logic-section]{subfiles}

\begin{document}

\olfileid{fol}{syn}{fvs}

\olsection{స్వేచ్ఛా \printtoken{P}{variable}, \printtoken{P}{sentence}}

\begin{defn}[\tetoken{ఒక చరం}{!!a{variable}} యొక్క స్వేచ్ఛా ఘటనలు]
\ollabel{defn:free-occ}
\tetoken{ఒక సూత్రంలోని}{!!a{formula}} \tetoken{ఒక చరం}{!!a{variable}} యొక్క \emph{స్వేచ్ఛా} ఘటనలను కింది
విధంగా ఆగమనాత్మకంగా నిర్వచిస్తాం:
\begin{enumerate}
\item \indcase*{!A}{$!A$ పరమాణువు}{ $\indfrm$లోని అన్ని \tetoken{చరం}{!!{variable}}
  ఘటనలు స్వేచ్ఛావే.}

\tagitem{prvNot}{\indcase{!A}{\lnot !B}{$\indfrm$లోని స్వేచ్ఛా
  \tetoken{చరం}{!!{variable}} ఘటనలు $!B$లో ఉన్నవే.}}{}

\item \indcase{!A}{(!B \ast !C)}{$\indfrm$లోని స్వేచ్ఛా
  \tetoken{చరం}{!!{variable}} ఘటనలు $!B$లోనివి, $!C$లోనివి.}

\tagitem{prvAll}{\indcase{!A}{\lforall[x][!B]}{$\indfrm$లోని స్వేచ్ఛా
  \tetoken{చరం}{!!{variable}} ఘటనలు $!B$లో ఉన్న వాటిలో~$x$ ఘటనలు మినహా మిగిలినవి.}}{}

\tagitem{prvEx}{\indcase{!A}{\lexists[x][!B]}{$\indfrm$లోని స్వేచ్ఛా
  \tetoken{చరం}{!!{variable}} ఘటనలు $!B$లో ఉన్న వాటిలో~$x$ ఘటనలు మినహా మిగిలినవి.}}{}
\end{enumerate}
\end{defn}

\begin{defn}[బద్ధ చరాలు]
ఒక సూత్రం~$!A$లోని \tetoken{ఒక చరం}{!!a{variable}} ఘటన స్వేచ్ఛగా లేకపోతే అది
\emph{బద్ధమైనది}.
\end{defn}

\begin{prob}
\olref[fol][syn][fvs]{defn:free-occ}ను అనుసరించే విధంగా బద్ధ
చర ఘటనలకు ఆగమనాత్మక నిర్వచనం ఇవ్వండి.
\end{prob}

\begin{defn}[పరిధి]
\iftag{prvAll}{ఒక సూత్రం~$!A$లో $\lforall[x][!B]$ ఒక ఉపసూత్ర ఘటన
  అయితే, $!A$లో దానికి అనుగుణంగా ఉన్న~$!B$ ఘటనను $\lforall[x]$కు
  అనుగుణమైన ఘటన యొక్క \emph{పరిధి} అంటాం. \iftag{prvEx}{అదే విధంగా
  $\lexists[x]$కూ.}{}}{ఒక సూత్రం~$!A$లో $\lexists[x][!B]$ ఒక
  ఉపసూత్ర ఘటన అయితే, $!A$లో దానికి అనుగుణంగా ఉన్న~$!B$ ఘటనను
  $\lexists[x]$కు అనుగుణమైన ఘటన యొక్క \emph{పరిధి} అంటాం.}

$!A$లో ఒక పరిమాణీకరణిక ఘటన
\iftag{prvAll}{$\lforall[x]$\iftag{prvEx}{ లేదా
    $\lexists[x]$}{}}{$\lexists[x]$}కు~$!B$ పరిధి అయితే, $!B$లోని $x$
స్వేచ్ఛా ఘటనలు
\iftag{prvAll}{$\lforall[x][!B]$లో\iftag{prvEx}{ మరియు
    $\lexists[x][!B]$లో}{}}{$\lexists[x][!B]$లో బద్ధమవుతాయి. ఈ
ఘటనలను పేర్కొన్న పరిమాణీకరణిక ఘటన \emph{బంధించింది} అని అంటాం.
\end{defn}

\begin{ex}
కింది సూత్రాన్ని పరిశీలించండి:
\[
\lexists[\Obj v_0][\underbrace{\Atom{\Obj A^2_0}{\Obj v_0,\Obj v_1}}_{!B}]
\]
$!B$, $\lexists[\Obj v_0]$ యొక్క పరిధిని సూచిస్తుంది.
పరిమాణీకరణికం $!B$లోని $\Obj v_0$ ఘటనను బంధిస్తుంది; కానీ
$\Obj v_1$ ఘటనను బంధించదు. అందువల్ల ఈ సందర్భంలో $\Obj v_1$ స్వేచ్ఛా
చరం.


ఇది మరింత సంక్లిష్టమైన \tetoken{సూత్రం}{!!{formula}}~$!A$లో ఎలా పనిచేస్తుందో ఇప్పుడు
చూడవచ్చు:
\[
\lforall[\Obj v_0][\underbrace{(\Atom{\Obj A^1_0}{\Obj v_0} \lif
    \Atom{\Obj A^2_0}{\Obj v_0, \Obj v_1})}_{!B}] \lif \lexists[\Obj
  v_1][\underbrace{(\Atom{\Obj A^2_1}{\Obj v_0, \Obj v_1} \lor \lforall[\Obj v_0][\overbrace{\lnot \Atom{\Obj A^1_1}{\Obj v_0}}^{!D}])}_{!C}]
\]
$!B$, మొదటి $\lforall[\Obj v_0]$ పరిధి; $!C$, $\lexists[\Obj v_1]$
పరిధి; $!D$, రెండవ $\lforall[\Obj v_0]$ పరిధి. మొదటి
$\lforall[\Obj v_0]$,~$!B$లోని $\Obj v_0$ ఘటనలను బంధిస్తుంది;
$\lexists[\Obj v_1]$,~$!C$లోని $\Obj v_1$ ఘటనను బంధిస్తుంది; రెండవ
$\lforall[\Obj v_0]$,~$!D$లోని $\Obj v_0$ ఘటనను బంధిస్తుంది. $!A$లోని
మొదటి $\Obj v_1$ ఘటన, నాలుగవ $\Obj v_0$ ఘటనలు స్వేచ్ఛావి. చివరి
$\Obj v_0$ ఘటన $!D$లో స్వేచ్ఛగా ఉన్నా, $!C$,~$!A$లలో బద్ధంగా ఉంది.
\end{ex}

\begin{defn}[వాక్యం]
\tetoken{ఒక సూత్రం}{!!^a{formula}}~$!A$లో \tetoken{చరాలు}{!!{variable}s} యొక్క స్వేచ్ఛా ఘటనలు ఏవీ లేకపోతే,
అప్పుడు మాత్రమే అది \article{\tetoken{వాక్యం}{!!{sentence}}}.
\end{defn}

% add examples!

\end{document}
